\section{Introduction}
\label{sec:intro}

\subsection{The CHESS mission}

CHESS (CubeSat for High-Energy Spectroscopy) is a CubeSat mission in
development at the EPFL Spacecraft Team. The satellite's payload is a
gamma-ray spectrometer designed to observe high-energy astrophysical
sources from low Earth orbit. CHESS is a student-led project: each
subsystem --- structure, attitude, thermal, communications, on-board
computer, and electrical power --- is owned by a team of student
engineers, with regular reviews against a master mission document
that defines requirements, modes, and interfaces.

The work described in this report concerns the
\textit{Electrical Power System} (EPS): the spacecraft subsystem that
generates, stores, and distributes electrical energy. The EPS is one
of the few subsystems that must operate before any other on the
satellite, because every other subsystem --- including the On-Board
Computer (OBC) that nominally supervises the spacecraft --- needs the
EPS to provide power before it can boot.

\subsection{The role of the EPS}

The CHESS EPS is built on three functional blocks:

\begin{description}
\item[Power Conditioning Unit (PCU)] Captures energy from the solar
array, regulates the charging current into the battery, and protects
the array--battery interface. The PCU contains the buck converter,
the maximum-power-point-tracking algorithm, and the four-mode
operating logic that governs how energy is converted in different
spacecraft conditions.

\item[Power Distribution Unit (PDU)] Provides regulated 3.3\,V, 5\,V,
and battery-bus rails to the rest of the satellite. The PDU contains
eFuses that the EPS can enable or disable through digital pins,
giving the satellite software-controlled load shedding for
contingency operations.

\item[EPS microcontroller (MCU)] Reads sensors, runs the control and
state-machine logic, commands the PCU and PDU actuators, reports
telemetry to the OBC, and falls back to autonomous behaviour if the
OBC stops responding. The MCU is a Microchip SAMD21 --- a 32-bit
Cortex-M0+ core with 128\,kB flash, 16\,kB RAM, and no external
operating system. All firmware described in this report runs on this
chip.
\end{description}

\subsection{Scope of this semester project}

The boundary of this semester project is the firmware that runs on
the MCU. The PCB hardware was designed by an earlier generation of
the team; component placement, schematic, and bill-of-materials are
fixed and were available throughout the work.

Inside that boundary, the project covered:

\begin{itemize}
\item designing a clean, layered firmware architecture from the
existing working but ad-hoc code base;
\item implementing the PCU operating-mode state machine in pure C,
without hardware dependencies, so that it could be tested on a
laptop before any board was involved;
\item implementing the Incremental Conductance MPPT algorithm with
the same separation between pure logic and hardware;
\item implementing the complementary 300\,kHz PWM driver with
hardware dead-time for the EPC2152 GaN half-bridge that drives the
buck converter;
\item implementing the CHIPS framing protocol used by every CHESS
subsystem to talk to the OBC;
\item building a hardware-in-the-loop demonstration harness centred
on a local Python web application and an ESP32 serial bridge, so
that the firmware running on the real chip can be commanded,
inspected, and exercised end-to-end from a laptop;
\item bringing up the real PCU testing board V4.1, including the
external programmer (Atmel-ICE through a Tag-Connect cable), the
buck-stage PWM scope verification, and a first round of injected
sensor-value tests;
\item planning, but not yet executing, a structured bench programme
to physically verify every firmware capability against known
electrical stimuli.
\end{itemize}

Out of scope for this semester were the physical sensor drivers for
the digital current monitors (INA226) and the SPI battery-temperature
sensor (whose part number is not yet specified in the mission
document), and the eventual integration of every subsystem into a
single super-loop running flight-grade per-mode logic. Both are
described in the future-work chapter as the next clearly-defined
steps.

\subsection{What ``done'' means for this report}

The semester deliverable is not a flight-ready EPS: the mission's
operational thresholds are not yet finalised, the battery part has
not been selected, and the spacecraft thermal model is still in
development. What the project does deliver is a firmware architecture
that can already be exercised end-to-end against real hardware, plus
the test infrastructure to validate every claim made about it.

Concretely, ``done'' for this report means:

\begin{enumerate}
\item every firmware capability listed in
Section~\ref{sec:intro}.3 is implemented and compiles cleanly under
the JPL/MISRA-derived coding rules of the project;
\item the firmware has been flashed and exercised on the real PCU
V4.1 mainboard, with the OBC link (an ESP32 acting as a stand-in)
verified by valid CHIPS responses;
\item the four PCU operating modes and the safe-mode entry triggers
have been exercised against injected sensor values, with each
expected mode reached;
\item the MPPT algorithm has been exercised against a software
solar-panel model on the ESP32, with closed-loop duty-cycle
convergence observed;
\item a structured bench-test plan exists with primary-source
probe-point references, so that any team member can take the project
forward without re-discovering the same PCB facts.
\end{enumerate}

\subsection{Report structure}

The remainder of the report is organised as follows.
\textbf{Chapter~\ref{sec:background}} introduces the mission
requirements, the electrical concepts (MPPT, buck converters,
complementary PWM, eFuses) needed to follow the rest of the report,
and the SAMD21 platform.
\textbf{Chapter~\ref{sec:architecture}} describes the firmware
architecture and walks through each major module.
\textbf{Chapter~\ref{sec:hil}} describes the hardware-in-the-loop test
harness --- the web app, the ESP32 bridge, and the demo pages --- and
the bench-test programme planned to physically verify each firmware
capability.
\textbf{Chapter~\ref{sec:results}} summarises the verified
behaviour, lists the gaps between the simulator-grade
implementation and a flight-grade EPS, and proposes a staged path
between them.
